\section{Supporting Results for Shared Laws}
\label{app:sharing}

\subsection{Own-arm-law limits and conditional stopping}

\begin{proof}[Proof of Proposition~\ref{prop:no-local}]
Let $q_0=0$. Arm $A$ has one collectable reward $R\sim\operatorname{Bernoulli}(1/2)$ at cost
$0.4$. Arm $B$ has a noncollectable signal $U\sim\operatorname{Bernoulli}(1/2)$ and a
collectable reward identically zero; each of its stages costs $0.01$. In the first instance
$U$ and $R$ are independent; in the second $U=R$. The separate arm laws are identical at the
root, with $G_A=0.2$ and $G_B=-0.02$ in both instances.
In the independent instance, starting $A$ has value $0.5-0.4=0.10$. Starting $B$ costs $0.01$
before the same optimal $A$ decision, giving $0.09$; its zero terminal reward is never worth
buying. In the shared instance, after observing $U=h$ the only potentially profitable action
is to inspect $A$, with net return $h-0.4$. Thus
$Q^\star(B)=-0.01+\E[\max\{0,R-0.4\}]=0.29$ exactly. Starting $A$ still has value $0.10$:
after its reward is known neither stage of $B$ can improve the incumbent or supply useful
information. The value $G_B=-0.02$ also follows from
$\Phi_B(r)=\max\{-0.01,(0-r)^+-0.02\}$ for a run restricted to $B$.
Stopping gives zero. Thus the unique optimal first actions are $A$ and $B$, respectively.
A rule with the stated inputs must choose the same action distribution in both, but optimality
requires probability one on different actions.
\end{proof}

For the general operational planning law, write
$\mathcal R=\{S:\max_{i\in\cI(S)}G_i(S)\le q(S)\}$,
with $\max\varnothing=-\infty$, for the arm-local index stopping region.
\begin{proposition}[Continuation and conditional stopping]\label{prop:stop-any}
Under Assumption~\ref{assm:standing}, at every $S\notin\mathcal R$,
\[
V^\star(S)\ge h(S)+\max_{i\in\cI(S)}\Phi_i^{(\nu_i)}(S;q(S))>h(S).
\]
If $\mathcal R$ is \emph{closed forward}, meaning that $F(S,i,Y)\in\mathcal R$ almost surely
for every $S\in\mathcal R$ and available $i$, then $V^\star(S)=h(S)$ on $\mathcal R$.
Under this additional hypothesis, $\mathcal R$ is exactly the optimal stopping region.
\end{proposition}
\begin{proof}
An attaining single-arm rule exists by finite-horizon backward induction.
Executing it and then stopping is admissible, with value
$h(S)+\Phi_i^{(\nu_i)}(S;q(S))$, proving the first assertion.
For the second, induct on the total stages remaining. At a state in $\mathcal R$, closure and
the induction hypothesis give $V^\star(S')=h(S')$ after every possible first advance.
Therefore
\[
Q^\star(S,i)=h(S)+\E[(\Yc-q(S))^+\mid S,i]-c(S,i)
\le h(S)+\Phi_i^{(\nu_i)}(S;q(S))\le h(S).
\]
The Bellman equation proves stopping optimal; the exhausted-state base case is immediate.
\end{proof}

\begin{example}[A premature index stop]\label{ex:shared-stop}
In the construction above, change $A$'s cost to $0.6$. Its index is now $-0.1$, while $B$'s
remains $-0.02$. The arm-local index stops at value zero, but buying $U$ and inspecting $A$
only when $U=1$ has value $-0.01+\tfrac12(1-0.6)=0.19$. Forward closure fails: after that
signal equals one, $A$'s index is $0.4>q_0$.
\end{example}

Closure quantifies over every available advance in $\mathcal R$, including those
the index never takes. It is sufficient, not necessary. Independent arms have
the exact stopping rule without it: an arm's own observation may raise its index
after an unprofitable initial inspection. The positive-run assertion concerns a
specified single-arm rule, rather than the utility of repeatedly recomputing and
switching under a general shared law. We do not establish closure for general
shared surrogates, including correlated Gaussian-process models.

\subsection{Proof of Proposition~\ref{prop:switching}}
\label{app:switching}

\begin{example}[The switching family]\label{ex:switching}
Let $\Theta\in\{0,1\}$ have prior $p=\Prob(\Theta=1)$ and let arms $A,B$, each with two
stages, be conditionally independent given $\Theta$. Arm $i$ has terminal reward
$R_i\in\{0,1\}$ with $\Prob(R_i=1\mid\Theta=0)=\alpha_i$ and $\Prob(R_i=1\mid\Theta=1)=1/5$,
where $\alpha_A=4/5$ and $\alpha_B=7/10$, and first reading $u_iS_i$, where $S_i\in\{0,1\}$
equals $R_i$ with probability $3/4$ given $R_i$, $u_A=2/5$ and $u_B=3/5$. First stages cost
$1/100$, terminal stages cost $c\in[19/100,1/5]$, and $q_0=0$. Arm $A$'s terminal reward is
more likely to be one when $\Theta=0$, and equally likely when $\Theta=1$;
arm $B$ has the larger intermediate score. Given the terminal rewards,
the signal errors are independent of $\Theta$ and of each other.
The thresholds are
\[
 \tau_1=\frac{15}{19},\qquad
 \tau_2(c)=\frac{132-420c}{50+35c},\qquad
 \tau_3(c)=\frac{288-780c}{250-535c}.
\]
\end{example}

\paragraph{A fixed-cost member.}
At $c=1/5$, the thresholds in Example~\ref{ex:switching} reduce to
$15/19$, $16/19$ and $12/13$. Terminal rewards are always collectable, and the
optimal first advances are
\begin{center}
\renewcommand{\arraystretch}{1.3}
\begin{tabular}{lcccc}
\toprule
Prior $p$ & $(0,\frac{15}{19})$ & $(\frac{15}{19},\frac{16}{19})$
 & $(\frac{16}{19},\frac{12}{13})$ & $(\frac{12}{13},1)$ \\
\midrule
$B$'s signal collectable & $A$ & $B$ & $A$ & $B$ \\
$B$'s signal not collectable & $A$ & $A$ & $A$ & $A$ \\
\bottomrule
\end{tabular}
\end{center}
Either choice of $A$'s intermediate collection rights gives the same rows.

We prove Proposition~\ref{prop:switching} for this family. Write $h=(h_A,h_B)$ for the observed history (each arm has an empty prefix, two
signal prefixes and four completed prefixes, so there are $49$ joint histories
$h$). Let $L_\theta(h)$ be the probability of $h$ given $\Theta=\theta$; every
one of the $16$ joint paths has positive probability in both latent states, so
no history reveals $\Theta$. For a fixed deterministic contingent policy, its
unnormalized payoff from $h$,
\[
 W^\pi_h(p,c)=\sum_\theta \Prob(\Theta=\theta)L_\theta(h)\,
 \E_\theta\Bigl[q(S_\tau)-\textstyle\sum_{t<\tau}c(S_t,i_t)\Bigm| h\Bigr],
\]
is affine in $p$ and affine in $c$, because the incumbent does not depend on
$c$ and each terminal stage contributes $-c$ times a path probability. The
unnormalized optimal value $W_h=\max_\pi W^\pi_h$ is therefore, for fixed $c$,
continuous and piecewise affine in $p$, and satisfies
\[
 W_h=\max\Bigl\{q(h)M_h,\;
 \max_{i}\Bigl[-c_i(h)M_h+\textstyle\sum_{y}W_{(h,i,y)}\Bigr]\Bigr\},
 \qquad M_h=(1-p)L_0(h)+pL_1(h),
\]
with $c_i(h)\in\{1/100,c\}$. Normalizing by $M_h>0$ gives the Bellman equation
for the posterior value, so the optimal first actions at the root are the
maximizers of the bracketed terms at $h=\varnothing$.

\emph{1. Certify all priors at a fixed cost.}
For rational $c$, compute the piecewise-affine Bellman recursion in rational
arithmetic. Partition $[0,1]$ at the union of breakpoints of each state's value,
action values and child values. On each interval, every expression is affine
in $p$: domination of all actions at both endpoints, with one action attaining
the value at both, proves the Bellman identity throughout. Backward induction
therefore certifies the optimum for every prior, rather than only sampled priors.
For example, with both intermediate readings collectable and $c=1/5$, the root
advance values with optimal continuation satisfy
\[
 Q_A(p)-Q_B(p)=\frac{15-19p}{2000},
 \qquad p\in[9/13,63/76].
\]
Their equality gives $\tau_1=15/19$ and the first change from $A$ to $B$;
the certificate also checks the winning advance against stopping.

\emph{2. Certify open intervals of costs.}
Run the same recursion and Bellman certificate over $\mathbb Q(c)$.
Every sort, breakpoint test and certificate comparison depends on the sign of
a rational-function difference $d$. Identically zero differences are exact
identities; for each nonzero $d$, record its sign at a generic rational sample
cost. If none has a zero or pole on an open interval $I$, every comparison
keeps its outcome there and evaluation commutes with the arithmetic.
Thus every $c\in I$ follows the same computational branch and inherits the
all-prior Bellman certificate. Exact Sturm counts on numerators and
denominators establish this sign invariance.

\emph{3. Cover exceptional costs and boundaries.}
On $(19/100,1/5)$, the only interior failure of sign invariance has numerator
proportional to $21000c^2-12035c+1536$, whose root is
$c^\ast=(2407-\sqrt{632689})/8400\approx0.19186$.
Here two partition points used by the computation coincide. This changes the
computation's partition, while the symbolic value and action functions it returns
are identical on the two sides. On $(19/100,c^\ast)$ and $(c^\ast,1/5)$,
all recorded nonzero differences keep their signs and the output gives the
thresholds in Example~\ref{ex:switching}. Fixed-cost certificates cover
$c=19/100$, $c=1/5$ in rational arithmetic and $c=c^\ast$ in
$\mathbb Q(\sqrt{632689})$.
With only $B$'s intermediate reading collectable, the same split and endpoint
checks give the same regions. With terminal-only collection or only $A$'s
intermediate reading collectable, no split is needed: $A$ is uniquely optimal
throughout $(0,1)$ for every cost in the closed interval.

Exact root counts also establish positive threshold denominators and
$0<\tau_1<\tau_2(c)<\tau_3(c)<1$ throughout the cost interval.
For uniqueness, within each open first-action region the winning advance strictly
exceeds stopping and the other advance at every contained breakpoint and midpoint;
nonnegative endpoint differences and positive midpoint differences of affine
functions give strict inequalities in the open piece. The advances tie at
$p=\tau_k(c)$. The verification script and output are included with the code.

\subsection{Signed correction and the regret bound}
\label{app:sharing-robustness}

Throughout this subsection impose the coherent-law and arm-local-cost restriction
of Proposition~\ref{prop:sharing-robustness}. Let
$\mathsf Q_S=\bigotimes_j \mathsf P_{j,S}$ and define
\[
 \mathcal V_\otimes(S)=\E_{\mathsf Q_S}\max\{q(S),Z_1^S,\ldots,Z_N^S\}.
\]
This lies in $[y_{\min},y_{\max}]$ and is the independent-surrogate optimum,
rather than the true shared optimum. Use the continuation utility
$U_\pi=V_\pi+K$ defined in Section~\ref{sec:setting}.

After forcing arm $i$ and seeing $Y=y$, write $S'=F(S,i,y)$. Define
$\mathcal V_i^{\rm fr}(S,y)$ using the new incumbent and the advanced arm's
conditional remaining law, but retain $\mathsf P_{j,S}$ for $j\ne i$. Coherence ensures
the advanced law equals its actual marginal at $S'$; only untouched laws differ.
Put
\begin{align}
 d_0(S)&=\mathcal V_\otimes(S)-q(S),\nonumber\\
 d_i(S)&=\mathcal V_\otimes(S)+c(S,i)-\E[\mathcal V_i^{\rm fr}(S,Y)\mid S,i],\nonumber\\
 \Gamma_i(S)&=\E[\mathcal V_\otimes(S')-\mathcal V_i^{\rm fr}(S,Y)\mid S,i].
 \label{eq:sharing-gamma}
\end{align}
Here $d_0$ is the independent surrogate's stop loss, $d_i$ its loss from
advancing $i$ under the frozen update, and $\Gamma_i$ the signed change in
the surrogate potential from updating untouched arms. Both signs of $\Gamma_i$
are possible; the proof controls its magnitude.
The losses $d_0,d_i$ are nonnegative by Theorem~\ref{thm:loss} applied to
$\mathsf Q_S$. In particular, with
$O_i=\max\{q,\max_{j\ne i}Z_j^S\}$ independent of arm $i$ under $\mathsf Q_S$,
\[
 d_i=c(S,i)-\E_{\mathsf Q_S}[(W_i^S-\max\{O_i,G_i(S)\})^+].
\]
For finite $G_i$ this is the capped-overshoot expression of Theorem~\ref{thm:loss};
for $G_i=-\infty$ it equals $c(S,i)$.

\begin{proof}[Proof of Proposition~\ref{prop:sharing-robustness}]
The true one-step residual is the signed update correction less the
nonnegative surrogate action loss:
\begin{equation}
 -c(S,i)+\E_{\mathsf P}[\mathcal V_\otimes(S')\mid S,i]-\mathcal V_\otimes(S)
       =\Gamma_i(S)-d_i(S).
 \label{eq:sharing-residual}
\end{equation}
Summing conditional expectations with indicators before the stopping time over
the deterministic range $0,\ldots,H-1$ gives, for every policy,
\begin{equation}
 U_\pi(S_0)=\mathcal V_\otimes(S_0)
 +\E_\pi\sum_{b<\tau_\pi}(\Gamma_{i_b}-d_{i_b})(S_b)
 -\E_\pi d_0(S_{\tau_\pi}).
 \label{eq:sharing-identity}
\end{equation}
All sums are finite and bounded. At each chosen index advance the surrogate
gap $d_i$ is zero; at its stop $d_0=0$. Consequently
\begin{equation}
 U_{\pi_{\rm IX}}(S_0)=\mathcal V_\otimes(S_0)
       +\E_{\pi_{\rm IX}}\sum_{b<\tau_{\rm IX}}\Gamma_{i_b}(S_b).
 \label{eq:sharing-index-identity}
\end{equation}
Subtracting from~\eqref{eq:sharing-identity} for a true optimal policy
$\pi^\star$, dropping its nonnegative stop gap, and taking positive parts yields
\begin{equation}
 U^\star-U_{\pi_{\rm IX}}
 \le\E_{\pi^\star}\sum_{b<\tau_\star}(\Gamma_{i_b}-d_{i_b})^+(S_b)
   +\E_{\pi_{\rm IX}}\sum_{b<\tau_{\rm IX}}\Gamma_{i_b}^-(S_b),
 \label{eq:sharing-refined}
\end{equation}
where $x^-=(-x)^+$. The first sum allows positive corrections along the true
optimal trajectory; the second allows negative corrections along the index
trajectory. This is why both trajectories enter the comparison.

Fix $S,i,y$. Replace the frozen untouched-arm capped distributions by their
updated distributions one arm at a time. For each replacement, condition on
all other independent capped variables. Their maximum with the new incumbent
is an outside option in $[q(S'),y_{\max}]$, so the absolute value change is at
most the corresponding supremum defining $\varepsilon_i(S)$. The advanced
arm's law is common to both products. Triangle inequality and expectation give
$|\Gamma_i(S)|\le\varepsilon_i(S)$. Each summand
in~\eqref{eq:sharing-refined} is at most $\bar\varepsilon$ and each trajectory
has at most $H$ advances. The two contributions are therefore at most
$H\bar\varepsilon$ each, giving $2H\bar\varepsilon$; removing the same sunk
cost from both values does not change regret.
\end{proof}

This is the finite-horizon approximate-planning argument. In the formulation of
\citet[Theorem~9 and Remark~11]{subramanian2022ais}, use the current marginal
laws and own histories as the abstract state, frozen updates as the approximate
transition, and incremental reward $q(S')-q(S)-c(S,i)$. The approximate value
is $\mathcal V_\otimes-q$ and its expected next-value discrepancy is exactly
$\Gamma_i$. The envelope above specializes that discrepancy using local
stop-loss values. The local values are $1$-Lipschitz in the outside option,
so their supremum is measurable by using a countable dense set and the
endpoints. Replacing $Z_j$ by $\max\{y_{\min},Z_j\}$ handles exhausted arms
without infinite arithmetic.

\paragraph{A path-total-variation envelope.}
Let $D=y_{\max}-y_{\min}$ and put $B_j(S)=D+\nu_j(S)c_{\max}$ if a
collectable stage remains, and zero otherwise. With
$\operatorname{TV}(\mu,\nu)=\sup_A|\mu(A)-\nu(A)|$,
\begin{equation}
 \varepsilon_i(S)\le\varepsilon_i^{\rm TV}(S)
 =\E\!\left[\sum_{j\ne i} B_j(S)
   \operatorname{TV}(\mathsf P_{j,S'},\mathsf P_{j,S})\,\middle|\,S,i\right].
 \label{eq:sharing-tv}
\end{equation}
With densities, our TV is half the $L^1$ distance; the convention in
\citet[Section~3.1]{subramanian2022ais} is the full $L^1$ distance.
The constants here use our stated convention throughout.
Indeed, each fixed local stopping rule, including immediate stopping, has path
payoff $\max\{r,M_\sigma\}-\kappa_\sigma$ in
$[y_{\min}-\nu_jc_{\max},y_{\max}]$. The same rule class is feasible under
both marginal path laws because costs are arm-local. TV bounds each expectation
difference by its payoff oscillation times TV, and taking the supremum over
rules preserves the bound. Without a collectable stage, immediate stopping is
optimal under both laws. This proves~\eqref{eq:sharing-tv}.

\begin{proof}[Proof of Corollary~\ref{cor:component}]
Conditional independence given $\Theta$ makes arm $j$'s remaining-path law at
$S$ the mixture $\int\mathsf P_j^\vartheta(\cdot\mid s_j)\,\mu_S(d\vartheta)$; a reading on
arm $i$ changes the mixing measure to $\mu_{S'}$ and leaves the components
unchanged. Write $\mu_{S'}-\mu_S=\xi^+-\xi^-$ by the Jordan decomposition, with
common mass $t=\operatorname{TV}(\mu_{S'},\mu_S)$. If $t=0$ the untouched
marginal laws agree; otherwise
$\mathsf P_{j,S'}-\mathsf P_{j,S}=t\bigl(\int\mathsf P_j^\vartheta\,d\xi^+/t-\int\mathsf P_j^\vartheta\,d\xi^-/t\bigr)$,
and the total variation between two mixtures of the components is at most the
largest total variation between components, so
$\operatorname{TV}(\mathsf P_{j,S'},\mathsf P_{j,S})\le\Lambda_j(S)\,t$.
Next, let $J$ be the joint law of $(\Theta,Y)$ given $S,i$, let $\beta$ be
its $Y$ marginal, and put $M=\mu_S\otimes\beta$. The finite measure $J+M$
dominates both $J$ and $M$. Integrating the absolute difference of their
densities with respect to $J+M$, first conditional on $Y$ and then conditional
on $\Theta$, gives the total-variation disintegration identity
\[
 \E_Y\operatorname{TV}(\mu_S(\cdot\mid Y),\mu_S)
 =\operatorname{TV}(J,M)
 =\E_{\vartheta\sim\mu_S}
   \operatorname{TV}(\mathsf P_i^\vartheta(Y\in\cdot\mid s_i),\beta)
 =\eta_i(S).
\]
Regular conditional laws exist on the stated standard Borel spaces.
This argument applies also to singular joint laws, without assuming $J\ll M$.
Substituting into~\eqref{eq:sharing-tv} gives the envelope, and
Proposition~\ref{prop:sharing-robustness} the regret bound. The two-draw bound
is the convexity inequality
$\operatorname{TV}(\mathsf P^\vartheta,\int\mathsf P^{\vartheta'}\mu(d\vartheta'))\le\int\operatorname{TV}(\mathsf P^\vartheta,\mathsf P^{\vartheta'})\mu(d\vartheta')$.
The binary formulas are the case of two components, computed next, in which $\varepsilon_i^{\rm TV}$ equals the displayed product.
\end{proof}

If arm paths are independent conditional on binary $\Theta$, let
$p=\mathsf P(\Theta=1\mid S)$, and let $D_i^{\rm next}$ and $D_j^{\rm path}$ be the
TV separations of the next-reading and remaining-path laws, respectively,
between the two latent components. Then, for $0<p<1$,
\begin{equation}
 \varepsilon_i^{\rm TV}(S)
 =2p(1-p)D_i^{\rm next}(S)\sum_{j\ne i} B_j(S)D_j^{\rm path}(S).
 \label{eq:sharing-binary}
\end{equation}
Here $\Theta$ is an uncertain latent state under the fixed planning law,
distinct from its fixed surrogate parameters $\theta$.
To see this, set $p_y=\mathsf P(\Theta=1\mid S,Y_i=y)$. An untouched arm's component
laws are unchanged, so its mixture-TV shift is
$|p_y-p|D_j^{\rm path}$. Bayes' rule gives
$\E|p_Y-p|=2p(1-p)D_i^{\rm next}$ by taking absolute variation of
$(p_y-p)\mu(dy)=p(1-p)(\mu_1-\mu_0)(dy)$ for
$\mu=(1-p)\mu_0+p\mu_1$. At $p\in\{0,1\}$ the envelope is zero, without
defining null-component conditionals.

\paragraph{Concentration and predictive dispersion.}
Weak posterior concentration alone does not force $\eta_i$ to vanish.
For example, consider beliefs $\mu_n=\operatorname{Uniform}[0,1/n]$ on a common
component $\Theta$, with every arm's next reading equal to $\Theta$.
These degenerate paths are conditionally independent given $\Theta$ and have
bounded readings. Although $\mu_n$ converges weakly to a point mass at zero,
$\operatorname{TV}(\delta_\vartheta,\mu_n)=1$ for every $\vartheta$, so
$\eta_i=1$ for every $n$. Predictive total-variation regularity is needed to
deduce local dispersion decay from concentration in a continuous component
space. To deduce a vanishing regret bound, the envelope must additionally
vanish uniformly over feasible future states, as required in
Corollary~\ref{cor:component}.

\paragraph{Why future-state control is required.}
Let $(U,A,B)$ be uniform over $000,011,101,110$. Each pair is independent.
Inspecting noncollectable $U$ costs $1/100$; inspecting collectable $A$ or $B$
costs $3/5$ each, with $q_0=0$. Every root envelope is zero because one root
observation leaves each other marginal unchanged, and the index stops.
After $U=1$, $A$ and $B$ are opposite: inspect $A$, then $B$ only if $A=0$,
for value $1/10$. After $U=0$, stopping is optimal. Buying $U$ therefore has
value $1/25$, whereas starting $A$ or $B$ has value $-1/200$. Later states have
nonzero envelopes. Thus even pairwise independence of whole arm paths at the
root cannot replace the all-history condition.
